% संपूर्ण मराठी ग्रंथातील प्रकरण 53: चिन्हार्थमीमांसा
% हा संपादनयोग्य स्रोतखंड आहे. संकलनासाठी संपूर्ण स्रोतसंग्रहातील
% अक्षररूपे, पूर्वभाग, चिन्हव्याख्या आणि आकृती-स्रोत आवश्यक आहेत.
% मूळ: Open Logic Project; मजकूर आणि रूपांतर CC BY 4.0.
% यंत्रानुवाद, दुरुस्ती आणि तपासणी: OpenAI Codex —
% GPT-5.6 Sol आणि GPT-6 Sol, दोन्ही Ultra effort.
% दुरुस्ती-पुनरावलोकन: GPT-6.1 Sol, Ultra effort.
\chapter{चिन्हार्थमीमांसा}\label{int:sem:chap}
\begin{editorial}
  या प्रकरणात अंतःप्रज्ञावादी तर्कशास्त्राच्या
  चिन्हार्थमीमांसेच्या व्याख्या एकत्र दिल्या
  आहेत. सध्या केवळ क्रिप्के चिन्हार्थमीमांसा आणि
  सांस्थितिक चिन्हार्थमीमांसा यांचा समावेश आहे.
  प्रतिमानांमध्ये सूत्रे सत्य कशी ठरतात किंवा
  सूत्रांची वैधता कशी सिद्ध करतात, यांची उदाहरणे
  अद्याप दिलेली नाहीत.
\end{editorial}
\section{प्रस्तावना}\label{int:sem:int:sec}
चिन्हार्थमीमांसेशिवाय कोणत्याही तर्कशास्त्राचे
समाधानकारक वर्णन होत नाही; अंतःप्रज्ञावादी
तर्कशास्त्रही त्याला अपवाद नाही. अभिजात
तर्कशास्त्रात सत्य-मूल्यांकने वर आधारलेली
चिन्हार्थमीमांसा प्रमाणभूत मानली जाते. पण
अंतःप्रज्ञावादी तर्कशास्त्रासाठी अनेक स्पर्धक
चिन्हार्थमीमांसा आहेत. संयोजकांच्या अर्थाचे
अंतःप्रज्ञावादी दृष्टिकोनातून मान्य होईल असे
पूर्ण स्पष्टीकरण देण्याच्या बाबतीत त्यांपैकी
एकही पूर्णपणे समाधानकारक नाही.
मोडल तर्कशास्त्रातील चिन्हार्थमीमांसेसारखी,
संबंधी प्रतिमाने वर आधारलेली चिन्हार्थमीमांसा
बहुधा सर्वाधिक प्रचलित आहे. यात विधानीय चले
जगांशी संबद्ध केली जातात आणि जगांना प्राप्यता
संबंध जोडतो. हा संबंध नेहमीच अंशतः क्रम असतो;
म्हणजे तो परावर्ती, प्रतिसममित आणि संक्रमक असतो.
सहज कल्पना करण्यासाठी या जगांना ज्ञानावस्था किंवा
“पुराव्याची परिस्थिती” समजू शकता. $w'$ ही अवस्था
$w$ पासून प्राप्य असेल तेव्हा आणि केवळ तेव्हाच,
आपल्याला सध्या माहीत असलेल्या गोष्टींनुसार $w'$
ही ज्ञानाची संभाव्य (भावी) अवस्था असते; म्हणजे
$w$ मध्ये जे माहीत आहे त्याच्याशी ती सुसंगत असते.
एकदा एखादे विधान माहीत झाले की ते नंतर अज्ञात
होऊ शकत नाही: $\varphi$ हे~$w$ मध्ये माहीत असेल
आणि~$Rww'$ असेल, तर $\varphi$ हे~$w'$ मध्येही माहीत
असते. त्यामुळे प्राप्यता संबंधाच्या दृष्टीने
“ज्ञान” एकस्वनिक आहे.
येथे “माहीत आहे” याचा अर्थ रचनाशील पडताळणीच्या
अर्थाने घ्यायचा आहे: सूत्र त्या ज्ञानावस्थेत सिद्ध झालेले
किंवा पडताळलेले आहे. ही पुढील विभागातील अंतःप्रज्ञावादी
सत्यतेची अट आहे; मागील ज्ञानविषयक तर्कशास्त्रातील
साधारण $\Knows_a$ संकारकाचा तोच अर्थ नाही.
$\varphi\land \psi$ हे~$w$ मध्ये पडताळलेले असते तेव्हा
आणि केवळ तेव्हाच $\varphi$ आणि $\psi$ दोन्ही तेथे पडताळलेले
असतात. $\varphi\lor \psi$ हे~$w$ मध्ये पडताळलेले असते
तेव्हा आणि केवळ तेव्हाच त्यांपैकी किमान एक तेथे
पडताळलेले असते. म्हणून $\land$ आणि $\lor$ यांच्या
स्थानिक सत्यतेच्या अटी अभिजात तर्कशास्त्रातील अटींशी जुळतात.
परावर्तिता आणि एकस्वनिकता यांमुळे आधी पडताळलेली सूत्रे
भावी अवस्थांतही पडताळलेली राहतात; मात्र “सर्व ज्ञानविषयक
पर्यायांत सत्य” या साधारण ज्ञान-संकारकाच्या अटीवरून
विकल्पासाठी ही स्थानिक अट निष्पन्न होत नाही.
उदाहरणार्थ, \ref{aml:el:trw:fig:simple} मधील $w_1$
येथे $\Knows_a(q\lor\lnot q)$ सत्य आहे, पण
$\Knows_a q\lor\Knows_a\lnot q$ असत्य आहे:
$a$ ला प्राप्य असलेल्या $w_1$ आणि $w_2$ या अवस्थांत
$q$ चे सत्यतामूल्य भिन्न आहे.
सशर्त विधानाच्या सत्यतेच्या अटी मात्र अभिजात
तर्कशास्त्राहून वेगळ्या आहेत. $\varphi \lif \psi$
हे~$w$ मध्ये माहीत असते तेव्हा आणि केवळ तेव्हाच
$Rww'$ असलेली अशी कोणतीही~$w'$ अवस्था नसते
जिथे $\varphi$ माहीत आहे पण $\psi$ माहीत नाही.
ही अट, $w$ मध्ये $\varphi$ माहीत नाही किंवा
$\psi$ माहीत आहे, या अटीसारखी नाही. कारण
$w$ मध्ये $\varphi$ किंवा $\psi$ यांपैकी काहीही माहीत
नसेल, तरी $Rww'$ असलेली अशी भावी ज्ञानावस्था
$w'$ असू शकते; त्या~$w'$ मध्ये $\varphi$ माहीत होते
पण $\psi$ माहीत होत नाही.
ज्या कोणत्याही संभाव्य भावी ज्ञानावस्थेत~$\varphi$
माहीत होते अशी अवस्था नसेल, तरच आपल्याला
$\lnot \varphi$ माहीत असते. यामागील कल्पना अशी:
$\varphi$ जाणून घेणे शक्य असेल, तर एखाद्या
संभाव्य भावी ज्ञानावस्थेत~$\varphi$ माहीत होते.
पण~$\lfalse$ माहीत असणे शक्य नाही; म्हणून
त्या भावी ज्ञानावस्थेत आपल्याला $\varphi$ माहीत
असेल, पण~$\lfalse$ माहीत नसेल.
या अर्थलक्षणात वर्जित मध्याचे तत्त्व वैध ठरत नाही.
कारण अशी काही $\varphi$ सूत्रे आहेत की जी अद्याप
माहीत नाहीत, पण पुढे माहीत होऊ शकतात.
अशा सूत्र~$\varphi$ साठी $\varphi$ आणि~$\lnot \varphi$
दोन्ही माहीत नसतात; म्हणून $\varphi \lor \lnot \varphi$
माहीत नसते. पण उदाहरणार्थ, $\lnot(\varphi
\land \lnot \varphi)$ हे आपल्याला माहीत असते.
कारण $\varphi$ आणि~$\lnot \varphi$ दोन्ही माहीत असलेली
कोणतीही भावी अवस्था शक्य नसते; आणि हे
$\varphi$ किंवा~$\lnot \varphi$ माहीत आहे की नाही
यापासून स्वतंत्रपणे आपल्याला माहीत असते.
अंतःप्रज्ञावादी तर्कशास्त्रासाठी संबंधी प्रतिमाने
हीच एकमेव उपलब्ध चिन्हार्थमीमांसा नाही.
सांस्थितिक चिन्हार्थमीमांसा हा आणखी एक पर्याय आहे:
त्यात विधानांचे अर्थलक्षण सांस्थितिक अवकाशातील
विवृत संच म्हणून केले जाते आणि संयोजकांचे
अर्थलक्षण त्या संचांवरील क्रिया म्हणून केले जाते
(उदा., $\land$ ला छेद अनुरूप असतो).
\section{संबंधी प्रतिमान}\label{int:sem:rel:sec}
अंतःप्रज्ञावादी विधानीय तर्कशास्त्राची अचूक
चिन्हार्थमीमांसा द्यायची असेल, तर ज्याच्या
संदर्भात सूत्रांचे मूल्यांकन करता येईल
असे प्रतिमान कोणते, याची व्याख्या द्यावी लागते.
तिच्या आधारे वैधता आणि निष्पन्नता यांसारख्या
चिन्हार्थविषयक संकल्पनांचीही व्याख्या करता येते.
अशी एक चिन्हार्थमीमांसा संबंधी प्रतिमाने
देतात.
\begin{defn}
  अंतःप्रज्ञावादी विधानीय तर्कशास्त्रासाठी
  संबंधी प्रतिमान म्हणजे त्रिक
  $\mModel{M} = \tuple{W, R, V}$, जिथे
  \begin{enumerate}
  \item $W$ हा अरिक्त संच आहे,
  \item $R$ हा~$W$ वरील अंशतः क्रम आहे (म्हणजे
    परावर्ती, प्रतिसममित आणि संक्रमक द्विस्थानी संबंध), आणि
  \item $V$ हे प्रत्येक विधानीय चल~$p$ ला~$W$ चा उपसंच नेमणारे फलन आहे;
    तसेच
  \item $V$ हे $R$ च्या दृष्टीने एकस्वनिक आहे:
    $w \in V(p)$ आणि $Rww'$ असेल, तर
    $w' \in V(p)$.
  \end{enumerate}
\end{defn}
\begin{defn}\label{int:sem:rel:defn:true-at-w}
  $\varphi$ हे \emph{$\mModel{M}$ मध्ये $w$ येथे सत्य असणे},
  म्हणजे $\mSat{M}{\varphi}[w]$, ही संकल्पना आपण पुढीलप्रमाणे
  विगमनाने परिभाषित करतो:
  \begin{enumerate}
  \item \indcase{\varphi}{p}{$\mSat{M}{\indfrm}[w]$ तेव्हा
    आणि केवळ तेव्हाच $w \in V(p)$.}
  \item \indcase{\varphi}{\lfalse}{$\mSat{M}{\indfrm}[w]$ नाही}.
  \item \indcase{\varphi}{\lnot \psi}{$\mSat{M}{\indfrm}[w]$
    तेव्हा आणि केवळ तेव्हाच $Rww'$ असलेले आणि
    $\mSat{M}{\psi}[w']$ पूर्ण करणारे कोणतेही
    $w'$ नाही}.
  \item \indcase{\varphi}{\psi \land \chi}{$\mSat{M}{\indfrm}[w]$
    तेव्हा आणि केवळ तेव्हाच $\mSat{M}{\psi}[w]$
    आणि $\mSat{M}{\chi}[w]$}.
  \item \indcase{\varphi}{\psi \lor \chi}{$\mSat{M}{\indfrm}[w]$
    तेव्हा आणि केवळ तेव्हाच $\mSat{M}{\psi}[w]$
    किंवा $\mSat{M}{\chi}[w]$ (किंवा दोन्ही)}.
  \item \indcase{\varphi}{\psi \lif \chi}{$\mSat{M}{\indfrm}[w]$
    तेव्हा आणि केवळ तेव्हाच $Rww'$ असलेल्या प्रत्येक
    $w'$ साठी $\mSat{M}{\psi}[w']$ नाही किंवा
    $\mSat{M}{\chi}[w']$} (किंवा दोन्ही).
  \end{enumerate}
  $\mSat{M}{\varphi}[w]$ नसल्यास आपण
  $\mSat/{M}{\varphi}[w]$ लिहितो. $\Gamma$ हा
  सूत्रांचा संच असेल, तर $\mSat{M}{\Gamma}[w]$
  याचा अर्थ प्रत्येक $\psi \in \Gamma$ साठी
  $\mSat{M}{\psi}[w]$ असा आहे.
\end{defn}
\begin{prob}
  \ref{int:sem:rel:defn:true-at-w} नुसार
  $\mSat{M}{\lnot \varphi}[w]$ तेव्हा आणि केवळ तेव्हाच
  $\mSat{M}{\varphi \lif \lfalse}[w]$, हे दाखवा.
\end{prob}
\begin{prop}\label{int:sem:rel:prop:true-monotonic}
  जगातील सत्यता $R$ च्या दृष्टीने एकस्वनिक आहे:
  $\mSat{M}{\varphi}[w]$ आणि $Rww'$ असेल, तर
  $\mSat{M}{\varphi}[w']$.
\end{prop}
\begin{proof}
  सराव.
\end{proof}
\begin{prob}
  \ref{int:sem:rel:prop:true-monotonic} सिद्ध करा.
\end{prob}
\section{चिन्हार्थविषयक संकल्पना}\label{int:sem:sem:sec}
\begin{defn}
  प्रत्येक $w \in W$ साठी $\mSat{M}{\varphi}[w]$
  असेल तेव्हा आणि केवळ तेव्हाच $\varphi$ हे
  \emph{$\mModel{M} = \tuple{W,R,V}$ या प्रतिमानात
  सत्य आहे}, म्हणजे $\mSat{M}{\varphi}$, असे म्हणतो.
  $\varphi$ हे सर्व प्रतिमानांत सत्य असेल तेव्हा आणि
  केवळ तेव्हाच ते \emph{वैध}, म्हणजे
  $\Entails \varphi$, असते. सूत्रांच्या
  संच~$\Gamma$ पासून~$\varphi$ \emph{तार्किकरीत्या
  निष्पन्न होते}, म्हणजे $\Gamma \Entails \varphi$,
  असे तेव्हा आणि केवळ तेव्हाच म्हणतो जेव्हा
  प्रत्येक प्रतिमान~$\mModel{M}$ आणि
  $\mSat{M}{\Gamma}[w]$ असलेल्या प्रत्येक~$w$
  साठी $\mSat{M}{\varphi}[w]$ असते.
\end{defn}
\begin{prop}\label{int:sem:sem:prop:sat-entails}
  \begin{enumerate}
  \item\label{int:sem:sem:prop:sat-entails1}
    $\mSat{M}{\Gamma}[w]$ आणि
    $\Gamma \Entails \varphi$ असेल, तर
    $\mSat{M}{\varphi}[w]$.
  \item\label{int:sem:sem:prop:sat-entails2}
    $\mSat{M}{\Gamma}$ आणि $\Gamma
    \Entails \varphi$ असेल, तर $\mSat{M}{\varphi}$.
  \end{enumerate}
\end{prop}
\begin{proof}
  \begin{enumerate}
  \item $\mSat{M}{\Gamma}[w]$ आणि
    $\Gamma \Entails \varphi$ असे समजा.
    निष्पन्नतेच्या व्याख्येवरून लगेच
    $\mSat{M}{\varphi}[w]$ मिळते.
  \item $\mSat{M}{\Gamma}$ असे समजा.
    मग प्रत्येक $u \in W$ साठी
    $\mSat{M}{\Gamma}[u]$; विशेषतः
    $\mSat{M}{\Gamma}[w]$. त्यामुळे
    \ref{int:sem:sem:prop:sat-entails1} नुसार
    $\mSat{M}{\varphi}[w]$. ते जग स्वेच्छ असल्याने
    निष्कर्ष प्रतिमानात सर्वत्र सत्य आहे.
  \end{enumerate}
\end{proof}
\begin{defn}\label{int:sem:sem:defn:restrict}
  $\mModel{M}$ हे संबंधी प्रतिमान आणि $w \in W$
  असे समजा. $w$ पासून प्राप्य जगांपुरते
  $\mModel{M}$ चे \emph{मर्यादन}
  $\mModel{M}_w=\tuple{W_w, R_w, V_w}$
  पुढीलप्रमाणे मिळते:
  \begin{align*}
    W_w & = \Setabs{u \in W}{Rwu},\\
    R_w  & = R \cap (W_w)^2, \text{ आणि}\\
    V_w(p) & = V(p) \cap W_w.
  \end{align*}
\end{defn}
\begin{prop}\label{int:sem:sem:prop:restrict}
  $\mSat{M}{\varphi}[w]$ तेव्हा आणि केवळ तेव्हाच
  $\mSat{M_w}{\varphi}$.
\end{prop}
\begin{prob}
  \ref{int:sem:sem:prop:restrict} सिद्ध करा.
\end{prob}
\begin{prop}
  $\mSat{M}{\Gamma}$ असलेल्या प्रत्येक
  प्रतिमान~$\mModel{M}$ मध्ये
  $\mSat{M}{\varphi}$ असेल, तर
  $\Gamma \Entails \varphi$.
\end{prop}
\begin{proof}
  $\mSat{M}{\Gamma}[w]$ असे समजा.
  प्रत्येक~$\psi \in \Gamma$ ला
  \ref{int:sem:sem:prop:restrict} लावल्यास
  $\mSat{M_w}{\Gamma}$ मिळते. गृहीतकावरून
  $\mSat{M_w}{\varphi}$ मिळते.
  पुन्हा \ref{int:sem:sem:prop:restrict} लावल्यास
  $\mSat{M}{\varphi}[w]$ मिळते.
\end{proof}
\section{सांस्थितिक चिन्हार्थमीमांसा}\label{int:sem:top:sec}
अंतःप्रज्ञावादी तर्कशास्त्राला चिन्हार्थमीमांसा
देण्याचा आणखी एक मार्ग म्हणजे संस्थिति ही
गणिती संकल्पना वापरणे.
\begin{defn}
  $X$ हा संच असू द्या. $X$ वरील
  \emph{संस्थिति} म्हणजे खालील गुणधर्म पूर्ण
  करणारा संच $\Top{O}
  \subseteq \Pow{X}$. $\Top{O}$ च्या
  घटक ना त्या संस्थितीचे
  \emph{विवृत संच} म्हणतात. $\Top{O}$
  सहित~$X$ या संचाला \emph{सांस्थितिक अवकाश}
  म्हणतात.
  \begin{enumerate}
  \item रिकामा संच आणि संपूर्ण अवकाश विवृत आहेत:
    $\emptyset$, $X \in \Top{O}$.
  \item विवृत संच सांत छेदांखाली संवृत असतात:
    $U$, $V \in \Top{O}$ असेल, तर
    $U \cap V \in \Top{O}$.
  \item विवृत संच स्वेच्छ संयोगांखाली संवृत असतात:
    प्रत्येक $i \in I$ साठी $U_i \in
    \Top{O}$ असेल, तर $\bigcup \Setabs{U_i}{i \in
      I} \in \Top{O}$.
  \end{enumerate}
\end{defn}
विवृत संचांचा संग्रह संदर्भावरून कळत असेल,
तर सांस्थितिक अवकाश दर्शवण्यासाठी फक्त $X$
लिहू शकतो. मात्र $X$ ला विवृत संचांचा संग्रह
दिल्यानंतरच त्याला सांस्थितिक अवकाश म्हणता येते.
\begin{defn}
  अंतःप्रज्ञावादी विधानीय तर्कशास्त्राचे
  \emph{सांस्थितिक प्रतिमान} म्हणजे
  $\mModel{X} = \tuple{X, \Top{O}, V}$ हे
  त्रिक, जिथे $\Top{O}$ ही~$X$ वरील
  संस्थिति आहे आणि $V$ हे प्रत्येक विधानीय
  चलाला~$\Top{O}$ मधील विवृत संच नेमणारे
  फलन आहे.
  सांस्थितिक प्रतिमान~$\mModel{X}$ दिले असता
  $\Prop{X}{\varphi}$ ची विगमनाने पुढीलप्रमाणे
  व्याख्या करता येते:
  \begin{enumerate}
  \item $\Prop{X}{\lfalse} = \emptyset$
  \item $\Prop{X}{p} = V(p)$
  \item $\Prop{X}{\varphi \land \psi} = \Prop{X}{\varphi} \cap \Prop{X}{\psi}$
  \item $\Prop{X}{\varphi \lor \psi} = \Prop{X}{\varphi} \cup \Prop{X}{\psi}$
  \item $\Prop{X}{\varphi \lif \psi} = \Interior{(X \setminus \Prop{X}{\varphi}) \cup
    \Prop{X}{\psi}}$
  \end{enumerate}
  येथे $\Interior{V}$ हे $V \subseteq X$
  या संचाला त्याच्या \emph{अंतर्भागाकडे} नेणारे
  फलन आहे; म्हणजे त्या संचात सामावलेल्या सर्व
  विवृत संचांचा संयोग. दुसऱ्या शब्दांत,
  \[
  \Interior{V} = \bigcup \Setabs{U}{U \subseteq V \text{ आणि } U \in \Top{O}}.
  \]
\end{defn}

कोणत्याही संचाचा अंतर्भाग हा विवृत संच असतो,
कारण तो विवृत संचांचा संयोग असतो. म्हणून
$\Prop{X}{\varphi}$ हा नेहमी विवृत संच असतो.

सांस्थितिक चिन्हार्थमीमांसा अत्यंत अमूर्त
असली, तरी तिच्यामागची प्रेरणा समजून घेण्याचे
मार्ग आहेत. $X$ चे घटक, म्हणजे
“बिंदू”, ही अशी स्थळे आहेत की जिथे विधानांचे
मूल्यांकन करता येते, असे समजा. ज्या सर्व
बिंदूंवर $\varphi$ सत्य आहे त्यांचा संच म्हणजे~$\varphi$
ने व्यक्त केलेले विधान. बिंदूंचा प्रत्येक संच
संभाव्य विधान असतोच असे नाही; फक्त $\Top{O}$ चे
घटक तसे असतात. $\varphi \Entails \psi$
तेव्हा आणि केवळ तेव्हाच $\varphi$ सत्य असलेल्या
प्रत्येक बिंदूवर $\psi$ सत्य असते; म्हणजे
प्रत्येक~$X$ आणि त्यावरील प्रत्येक अनुमत मूल्यांकनासाठी $\Prop{X}{\varphi}
\subseteq \Prop{X}{\psi}$. असंगत विधान~$\lfalse$
कधीच सत्य नसते; म्हणून
$\Prop{X}{\lfalse} = \emptyset$.

अंतःप्रज्ञावादी तर्कशास्त्रात वैध असणारे
नियम टिकावेत, म्हणजे सर्व सांस्थितिक प्रतिमानांसाठी
$\varphi \Proves \psi$ तेव्हा आणि केवळ तेव्हाच
$\Prop{X}{\varphi} \subseteq \Prop{X}{\psi}$
असावे, यासाठी $\psi \land \chi$,
$\psi \lor \chi$ आणि $\psi \lif \chi$ यांनी
व्यक्त केलेली विधाने $\psi$ आणि~$\chi$
यांनी व्यक्त केलेल्या विधानांशी कशी संबंधित
असावीत? कोणत्याही $\varphi$ साठी
$\lfalse \Proves \varphi$ अपेक्षित आहे; कारण
प्रत्येक $U$ साठी $\emptyset
\subseteq U$ असल्याने ते पूर्ण होते.
$\psi \land \chi \Proves \psi$ असल्याने
$\Prop{X}{\psi \land \chi} \subseteq \Prop{X}{\psi}$
ही अट हवी; तसेच
$\Prop{X}{\psi \land \chi} \subseteq \Prop{X}{\chi}$.
$W \subseteq U$ आणि $W \subseteq V$
या अटी पूर्ण करणारा सर्वांत मोठा संच
$U \cap V$ आहे. उलट,
$\psi \Proves \psi \lor \chi$ आणि
$\chi \Proves \psi \lor \chi$, म्हणून
$\Prop{X}{\psi} \subseteq \Prop{X}{\psi \lor \chi}$
आणि $\Prop{X}{\chi} \subseteq \Prop{X}{\psi \lor \chi}$
या अटी हव्यात. $U \subseteq W$ आणि
$V \subseteq W$ असलेला सर्वांत लहान
संच~$W$ म्हणजे $U \cup V$.

$\lif$ साठीची व्याख्या अधिक सूक्ष्म आहे:
$\varphi \lif \psi$ हे, $\varphi$ शी संयोग केल्यावर
$\psi$ निष्पन्न होईल, असे सर्वांत दुर्बल
विधान व्यक्त करते. $\varphi \lif \psi$ चा
$\varphi$ शी संयोग केल्यावर~$\psi$ निष्पन्न होते,
हे $(\varphi \lif \psi)
\land \varphi \Proves \psi $ वरून स्पष्ट आहे.
म्हणून $\Prop{X}{\varphi \lif \psi}$ हा
$\Prop{X}{\varphi \lif \psi} \cap \Prop{X}{\varphi}
\subseteq \Prop{X}{\psi}$ ही अट पूर्ण करणारा
सर्वांत मोठा विवृत संच असला पाहिजे;
यावरून आपली व्याख्या मिळते.



